\section{March 20}

\subsection{Outline of proof of local CFT: Methods 2 and 3}

  \textbf{Continue:}
  The second method of proving the local CFT is to combine the Lubin-Tate method with cohomological methods.
  If we have proved the main theorem, then for all uniformizer $\pi$ of $K$ and all $n\geq 1$, there is an abelian extension $K_{\pi,n}/K$ with $\Nm(K_{\pi,n}) = (1+\frakm^n)\pi^\bbZ$.
  Then they have the following properties:
  \begin{enumerate}
    \item $[K_{\pi,n}:K] = (q-1)q^{n-1}$;
    \item $\forall n$, $\pi$ is a norm from $K_{\pi,n}$.
  \end{enumerate}
  Moreover, $K^\ab = K_\pi \cdot K^\unr$.
  In this method, we explicitly construct a totally ramified abelian extension satisfying the above properties.
  Let $K_\pi = \bigcup_{n\geq 1} K_{\pi,n}$.
  For all $m \geq 1$, let $K_m$ be the unique unramified extension of $K$ of degree $m$.
  We explicitly construct $\phi_\pi: K^\times \to \Gal(K_\pi \cdot K^\unr/K)$ such that 
  \begin{itemize}
    \item[(c)] $\phi_\pi(\pi)|_{K^\unr} = \Frob_K$;
    \item[(d)] $\phi_\pi(a)|_{K_{\pi,n} \cdot K^\unr} = \id$ for all $a \in (1+\frakm^n) \cdot \pi^{m\bbZ}$.
  \end{itemize}
  % $\phi_\pi(\pi)|_{K^\unr} = \Frob_K$ and $\phi_\pi(a)|_{K_{\pi,n} \cdot K^\unr} = \id$ for all $a \in (1+\frakm^n) \cdot (\pi^m)$.
  Moreover, we show that both $K_\pi \cdot K^\unr$ and $\phi_\pi$ are independent of the choice of $\pi$.
  Via cohomological methods, we can show the existence in \cref{thm:local_reciprocity_law}.
  Now we explain how these results implies the local CFT.

  Let $K' = K_\pi \cdot K^\unr \subset K^\ab$ and $\phi' = \phi_\pi: K^\times \to \Gal(K'/K)$.

  \begin{lemma}
    For all $a \in K^\times$, $\phi(a)|_{K'} = \phi'(a)$.
  \end{lemma}
  \begin{proof}
    For all uniformizer $\pi$, we have $\phi(\pi)|_{K_{\pi,n}} = \id$ as $\pi$ is a norm from $K_{\pi,n}$.
    And $\phi'(\pi)|_{K_{\pi,n}} = \id$ by (d).
    Both $\phi(\pi)$ and $\phi'(\pi)$ act as $\Frob_K$ on $K^\unr$.
    They must agree on $K' = K_\pi \cdot K^\unr$.
    Since $K^\times$ is generated by uniformizers and $\phi'$ is independent of the choice of $\pi$, we have $\phi(a)|_{K'} = \phi'(a)$ for all $a \in K^\times$.
    % \Yang{}
  \end{proof}

  \begin{theorem}[Local Kronecker-Weber]
    For all uniformizers $\pi$ of $K$, we have $K^\ab = K_\pi \cdot K^\unr$.
  \end{theorem}
  \begin{proof}
    Let $K_{n,m} = K_{\pi,n} \cdot K_m$ and $U_{n,m} = (1+\frakm^n) \cdot (\pi^m)$.
    % By \cref{thm1}, $\phi(a)|_{K_{n,m}} = \id$ for all $a \in U_{n,m}$.
    By our construction of $\phi_\pi$, $\phi_\pi(a)|_{K_{n,m}} = \id$ for all $a \in U_{n,m}$.
    Hence $\phi(a)|_{K_{n,m}} = \id$ for all $a \in U_{n,m}$.
    Hence $U_{n,m} \subset \Nm(K_{n,m}/K)$.
    We have 
    \[ [K^\times : U_{n,m}] = [U:1+m^n] \cdot [\pi^\bbZ : \pi^{m \bbZ}] = (q-1)q^{n-1}m = [K_{n,m}:K]. \] 
    Hence $U_{n,m} = \Nm(K_{n,m})$.
    Let $L/K$ be a finite abelian extension.
    Then $\Nm(L) \subset K^\times$ is an open subgroup.
    Then there exists $n,m$ such that $\Nm(L) \supset U_{n,m}$.
    Hence $L \subset K_{n,m} \subset K'$ by \cref{cor:correspondence_between_extensions_and_norm_groups}.
    Take union over all $L$, we have $K^\ab \subset K'$.
  \end{proof}

  Hence we only need to show \cref{thm:norm_groups_are_exactly_open_subgroups_of_finite_index} any open subgroup $H$ of $K^\times$ of finite index is a norm group.
  This is since $H \supset U_{n,m}$ for some $n,m$ and $U_{n,m}$ is a norm group.
  Then the assertion follows from \cref{cor:correspondence_between_extensions_and_norm_groups}.


% \subsection{}

  \paragraph{The third method of proving local CFT:}
  This method is to use cohomological methods and Hilbert symbols.
  Using cohomological methods, we can show the existence of local Artin map.
  Using Hilbert symbols, we can show \cref{thm:norm_groups_are_exactly_open_subgroups_of_finite_index}.



\subsection{Lubin-Tate: formal group laws}

  Let $A$ be a commutative ring with identity.
  Recall
  \[ A[[T]] = \left\{ \sum_{n=0}^\infty a_n T^n : a_n \in A \right\} \]
  is an algebra over $A$.
  For $f,g \in A[[T]]$ with $g(0) = 0$, we can define the composition $f \circ g$ by
  \[ (f \circ g)(T) = f(g(T)) = \sum_{n=0}^\infty a_n (g(T))^n. \]

  \begin{lemma}
    \begin{enumerate}
      \item If $f \in A[[T]]$ with $g,h \in T A[[T]]$, then $f \circ (g \circ h) = (f \circ g) \circ h$.
      \item Let $f \in T A[[T]]$. Then there exists $g \in T A[[T]]$ such that $f \circ g = T$ iff $a_1 \in A^\times$.
        In which case $g$ is unique and $g \circ f = T$. 
        We denote $g$ by $f^{-1}$.
    \end{enumerate}
  \end{lemma}

  Power series in serval variables can also be defined similarly.
  More over, if $f \in A[[T_1,\cdots,T_n]]$ and $g_1,\cdots,g_n \in A[[T_1,\cdots,T_m]]$ with $g_i(0) = 0$, then the composition $f \circ (g_1,\cdots,g_n) = f(g_1(T),\cdots,g_n(T))$ is well-defined.

  \begin{remark}
    Let $(A,\frakm)$ be a complete DVR.
    Then for all $f = \sum_{n=0}^\infty a_n T^n \in A[[T]]$ and $c \in \frakm$, we have $a_n c^n \to 0$ as $n \to \infty$ and then $f(c)$ is well-defined in $A$.
    Moreover, $f(c) \in \frakm$ if $f(0) \in \frakm$.
  \end{remark}

  \begin{slogan}
    A formal group law is a law of composition (without set) satisfying the group axioms.
  \end{slogan}

  \begin{definition}
    Let $A$ be a ring.
    A \emph{one-parameter commutative formal group law} is a power series $F \in A[[X,Y]]$ such that
    \begin{enumerate}
      \item $F(X,Y) \equiv X+Y + \epsilon_{\geq 2}$;
      \item $F(X,F(Y,Z)) = F(F(X,Y),Z)$;
      \item $\exists ! i_{F} \in T A[[T]]$ such that $F(T,i_F(T)) = 0$;
      \item $F(X,Y) = F(Y,X)$.
    \end{enumerate}
    Here $\epsilon_{\geq 2}$ is a power series in $X,Y$ with all monomials of degree $\geq 2$.
  \end{definition}

  \begin{remark}
    Condition (a) implies that $F(0) = 0$ and then (b) makes sense.
  \end{remark}

  \begin{remark}
    Taking $Y=Z = 0$ in (a) and (b), we have $F(X,0) = X + \epsilon_{\geq 2}$ and $F(F(X,0),0) = F(X,0)$.
    Set $f(X) = F(X,0)$. 
    Then $f(X)$ is invertible and $f = f^{\circ 2}$.
    Set $g(X) = f^{-1}(X)$. 
    Then $f = f\circ (f \circ g) = (f \circ f) \circ g = f \circ g = X$.
    Hence $F(X,0) = X$.
    Similarly, $F(0,Y) = Y$.
    This can be viewed as the existence of identity element $0$.
  \end{remark}

  Note that above conditions imply that $F(X,Y)$ is of the form
  \[ F(X,Y) = X + Y + \sum_{i,j \geq 1} a_{i,j} X^i Y^j. \]

  In this note, the term ``formal group law'' always means one-parameter commutative formal group law.

  Let $A = \calO_K$ with $K$ a non-archimedean local field and $F = \sum a_{ij} X^i Y^j$ a formal group law over $A$.
  For any $x,y \in \frakm_K$, we have $a_{ij} x^i y^j \to 0$ as $i+j \to \infty$ and then $F(x,y)$ is well-defined in $\frakm_K$.
  Denote by $x +_F y$ the value of $F(x,y)$.
  In this way, $F$ defines a group structure on $\frakm_K$.
  For any finite extension $L/K$, we can also define a group structure on $\frakm_L$ by $F$.
  The inclusion $\frakm_K \subset \frakm_L$ is a group homomorphism.

  \begin{example}
    Let $F(X,Y) = X + Y$.
    Then $x +_F y = x+y$ is the usual addition.
  \end{example}

  \begin{example}
    Let $F(X,Y) = X + Y + XY$.
    The map $a \mapsto a + 1$ is a group isomorphism from $(\frakm_K, +_F)$ to $(1+\frakm_K, \cdot)$.
  \end{example}

  \begin{example}
    Consider an elliptic curve $E$ over $K$.
    Expand the addition law of $E$ at the origin $O$ as a power series in $X,Y$.
    We get a formal group law $F_E$ over $K$.
  \end{example}

  \begin{definition}
    Let $F,G$ be two formal group laws over $A$.
    A \emph{homomorphism} from $F$ to $G$ is a power series $h \in T A[[T]]$ such that $h(F(X,Y)) = G(h(X),h(Y))$.
    When $h$ has an inverse, we say $h$ is an \emph{isomorphism} and $F$ and $G$ are \emph{isomorphic}.
  \end{definition}

  A homomorphism $h$ from $F$ to $G$ defines a group homomorphism from $(\frakm_L, +_F)$ to $(\frakm_L, +_G)$ for all finite extension $L/K$.

  % \begin{definition}
  %   Let $F = X + Y + XY$.
  % \end{definition}

  \begin{example}
    Let $F(X,Y) = (1+X)(1+Y) - 1$ and $f(T) = (1+T)^p - 1$.
    Then $f(F(X,Y)) = f(X) + f(Y) + f(X)f(Y) = F(f(X),f(Y))$.
    Hence $f$ is a homomorphism from $F$ to itself.
    Consider 
    \[ \begin{tikzcd}
      \frakm_K \arrow[r, "f"] \arrow[d, "+1"] & \frakm_K \arrow[d, "+1"] \\
      1+ \frakm_K \arrow[r, "a \mapsto a^p"] & 1+ \frakm_K.
    \end{tikzcd} \]
    Hence $f$ defines an homomorphism on $(\frakm_K, +_F) \cong (1+\frakm_K, \cdot)$ which is the $p$-th power map.
  \end{example}

  Let $G$ be a formal group law.
  For all $f,g \in T A[[T]]$, define
  \[ f+_G g \coloneqq G(f,g). \]
  Then $T A[[T]]$ is a commutative group under $+_G$ by axioms of formal group law.

  \begin{lemma}
    \begin{enumerate}
      \item For any formal group laws $F,G$, then set $\Hom(F,G)$ of homomorphisms from $F$ to $G$ becomes an abelian group under $+_G$.
      \item For any formal group laws $F$, $\End(F) = \Hom(F,F)$ is a ring under $+_F$ and composition (not necessarily commutative).
    \end{enumerate}
  \end{lemma}
  \begin{proof}
    Exercise.
  \end{proof}
